\documentclass[a5paper,10pt]{article}

% https://github.com/InDevRus/filippov-solutions

% Нумерация страниц
\pagenumbering{gobble}

% Настройка полей
\usepackage{geometry}
\geometry{left=1cm}
\geometry{right=1cm}
\geometry{top=1cm}
\geometry{bottom=1.5cm}

% Вставка таблицы
\usepackage{array}
\usepackage{tabularx}

% Инструменты выравнивания формул
\usepackage{amsmath}
\usepackage{mathtools}

% Диакритика в математике
\usepackage{amsfonts}

% Вычисления размеров на ходу
\usepackage{calc}

% Удобные записи производных
\usepackage[italic=true]{derivative}

% Настройки сносок
\usepackage{enotez}

% Длинные знаки векторов
\usepackage{anyfontsize}
\usepackage[b]{esvect}

% Вставка картинок
\usepackage{graphicx}

% Вставка красной строки
\usepackage{indentfirst}
\setlength{\parindent}{1em}

% Условия в командах
\usepackage{xifthen}

% Луа-скрипты
\usepackage{luacode}

% Настройка команд отдельно в математическом и текстовом режимах
\usepackage{mathcommand}

% Для повторения LaTeX кода
\usepackage{multido}

% Конфигурация отступов формул
\delimitershortfall-1sp
\usepackage{mleftright}
\mleftright

% Растягивание объектов
\usepackage{relsize}

% Обтекание текстом
\usepackage{wrapfig}
\usepackage{subcaption}
\usepackage{floatflt}

% Поддержка ввода кириллицы
\usepackage[english, russian]{babel}

% Настройка корневой позиции для компилятора
\usepackage{import}
\providecommand*{\ProjectRootPath}{..}

\import{\ProjectRootPath/preambles/macros/}{theorems}

\import{\ProjectRootPath/preambles/fonts}{font_setup}

% Знак градуса
\usepackage{siunitx}

\import{\ProjectRootPath/preambles/macros/}{operators}
\import{\ProjectRootPath/preambles/macros/}{redefinitions}

\import{\ProjectRootPath/preambles/fonts}{letters_setup}
\import{\ProjectRootPath/preambles/fonts/adjustments}{custom_superscripts}

\import{\ProjectRootPath/preambles/custom}{formatting}
\import{\ProjectRootPath/preambles/custom}{frames}
\import{\ProjectRootPath/preambles/custom}{list_setup}

% TODO: перевести команды на современный стиль объявления
\input{../preambles/plotting}

\begin{document}

$\ddot{x} = 4\dot{x} - 5x$.
$\tsys{& \dot{x}\br{t} = y \\& \dot{y}\br{t} = -3x + 4y}$;
$\dfrac {dy} {dx} = \dfrac {-3x + 4y} {y}$.

Особой точкой является $M\br{0; 0}$.
$\begin{pmatrix} 0 & 1 \\ -3 & 4 \end{pmatrix}$.
$\lambda^2 -4\lambda + 3 = 0$.
$\lambda_{1} = 1$, $\lambda_{2} = 3$.
Таким образом, $M\br{0; 0}$ -- <<узел>>, для написанной выше системы неустойчивый.
Для $\lambda_{1} = 1$ собственный вектор $\vv{v}_{1} = \cbr{1; 1}$ отвечает касательной $y = x$; для $\lambda_{2} = 3$ собственным вектором будет $\vv{v}_{2} = \cbr{1; 3}$, и он задаёт трансверсаль $y = 3x$.

Уравнения траекторий имеют вид $y - x = C\br{y - 3x}^3$ (в неё не входит только $y = 3x$), либо же в параметрической форме $\tsys{& x\br{\tau} = -\tau + C \tau^3 \\& y\br{\tau} = -\tau + 3C \tau^3}$.
	
\begin{floatingfigure}{0.4\textwidth}
    \begin{tikzpicture}
        \pgfplotsset{
            Graph/.style = {
                samples = 500,
                restrict y to domain = -5 : 5,
            },
            DirectionGraph/.style = {
                samples = 20,
                quiver = {scale arrows = 0.075},
                restrict y to domain = -5 : 5,
                -stealth,
            },
            DirectionGraph1/.style = {
                DirectionGraph,
                quiver = {
                    u = {dot_x(x, y) / sqrt(dot_x(x, y) ^ 2 + dot_y(x, y) ^ 2)},
                    v = {dot_y(x, y) / sqrt(dot_x(x, y) ^ 2 + dot_y(x, y) ^ 2)}
                }
            },
            DirectionGraph2/.style = {
                DirectionGraph,
                quiver = {
                    u = {-dot_x(x, y) / sqrt(dot_x(x, y) ^ 2 + dot_y(x, y) ^ 2)},
                    v = {-dot_y(x, y) / sqrt(dot_x(x, y) ^ 2 + dot_y(x, y) ^ 2)}
                }
            },
        }
        \begin{axis}[
            trig format plots = rad,
            width = 6.5cm,
            height = 10cm,
            ylabel=$\dot{x}$,
            ymin = -1.9,
            ymax = 1.9,
            xmin = -1.6,
            xmax = 1.6, 
            xtick = {-10, ..., 10},
            ytick = {-10, ..., 10},
            minor tick num = 3,
            declare function = {
            	f(\C, \x) = -x / 2 + \C * x^3;
            	g(\C, \x) = -x / 2 + 3 * \C * x^3;
            	dot_x(\x, \y) = y;
            	dot_y(\x, \y) = -3 * x + 4 * y;
            }]
            
            \addplot[Graph, dashed, domain = -2 : 2] 
            {x};
            \addplot[Graph, dotted, domain = -2 : 2] 
            {3 * x};
            
            \pgfmathsetmacro{\k}{1}
            \foreach \s in {-1, 1} {
            	\addplot[Graph, olive, domain = -5 : 5]
            	({f(\s * \k, x)}, {g(\s * \k, x)});
            }
            
            \pgfmathsetmacro{\k}{5}
            \foreach \s in {-1, 1} {
            	\addplot[Graph, blue, domain = -5 : 5]
            	({f(\s * \k, x)}, {g(\s * \k, x)});
            	\addplot[DirectionGraph1, blue, samples = 6, domain = -0.5 : -0.2]
            	({f(\s * \k, x)}, {g(\s * \k, x)});
            	\addplot[DirectionGraph1, blue, samples = 6, domain = 0.2 : 0.5]
            	({f(\s * \k, x)}, {g(\s * \k, x)});
            }
            
            \pgfmathsetmacro{\k}{0.1}
            \foreach \s in {-1, 1} {
            	\addplot[Graph, red, domain = -5 : 5]
            	({f(\s * \k, x)}, {g(\s * \k, x)});
            	\addplot[DirectionGraph1, red, samples = 6, domain = -2 : -0.6]
            	({f(\s * \k, x)}, {g(\s * \k, x)});
            	\addplot[DirectionGraph1, red, samples = 6, domain = 0.6 : 2]
            	({f(\s * \k, x)}, {g(\s * \k, x)});
            }
            
            \pgfmathsetmacro{\k}{0.025}
            \foreach \s in {-1, 1} {
            	\addplot[Graph, magenta, domain = -10 : 10]
            	({f(\s * \k, x)}, {g(\s * \k, x)});
            	\addplot[DirectionGraph1, magenta, samples = 9, domain = -4 : -1.5]
            	({f(\s * \k, x)}, {g(\s * \k, x)});
            	\addplot[DirectionGraph1, magenta, samples = 9, domain = 1.5 : 4]
            	({f(\s * \k, x)}, {g(\s * \k, x)});
            }
            
            \pgfmathsetmacro{\k}{0.01}
            \foreach \s in {-1, 1} {
            	\addplot[Graph, purple, domain = -10 : 10]
            	({f(\s * \k, x)}, {g(\s * \k, x)});
            	\addplot[DirectionGraph1, purple, samples = 9, domain = -6 : -1.5]
            	({f(\s * \k, x)}, {g(\s * \k, x)});
            	\addplot[DirectionGraph1, purple, samples = 9, domain = 1.5 : 6]
            	({f(\s * \k, x)}, {g(\s * \k, x)});
            }
        \end{axis}
    \end{tikzpicture}
\end{floatingfigure}
Решения $x\br{t}$ при $t \to +\infty$ стремятся к $+\infty$.

\end{document}
